% ============================================================================
% cluster_distribution_theory.tex
% AI-DRAFTED (Claude), 2026-08-13. Seam protocol: this whole file is one
% generated block. Source: agents/2026-08-13_claude_cluster_distribution_analysis.md
% and agents/2026-08-13-thoughts_cluster_distribution.md.
% Mechanism-candidate sentences (owner rewrites) are marked with %% MECHANISM.
% Compile: pdflatex + bibtex with cluster_lineage_refs.bib (plainnat).
% The target-journal class and bibliography style override this preamble.
% ============================================================================
\documentclass[11pt]{article}
\usepackage[margin=2.7cm]{geometry}
\usepackage{amsmath, amssymb, amsthm}
\usepackage{natbib}
\usepackage{hyperref}

\newtheorem{lemma}{Lemma}
\newtheorem{proposition}{Proposition}
\newtheorem{corollary}{Corollary}
\newtheorem{assumption}{Assumption}
\newtheorem{prediction}{Prediction}
\theoremstyle{remark}
\newtheorem{remark}{Remark}

\newcommand{\E}{\mathbb{E}}
\newcommand{\Prob}{\mathbb{P}}
\newcommand{\R}{\mathbb{R}}
\newcommand{\vech}{\operatorname{vech}}
\newcommand{\ARI}{\operatorname{ARI}}
\newcommand{\VI}{\operatorname{VI}}

\title{A distribution theory for rolling correlation-based clusters\\
\large Working draft for the cluster-lineage paper}
\author{Artur Sepp\thanks{Draft section prepared for the Quantitative Finance
manuscript. Reference implementation in the \texttt{factorlasso} package.}}
\date{2026-08-13}

\begin{document}
\maketitle

% ----------------------------------------------------------------------------
\section{Setup and notation}
% ----------------------------------------------------------------------------

We observe returns $r_t \in \R^d$ on $d$ assets at a fixed frequency. The
covariance estimator is the exponentially weighted moving average (EWMA) with
span $N$,
\begin{equation}
\widehat{\Sigma}_t \;=\; \lambda\, \widehat{\Sigma}_{t-1}
\;+\; (1-\lambda)\, r_t r_t^{\top},
\qquad \lambda \;=\; 1 - \frac{2}{N+1},
\label{eq:ewma}
\end{equation}
with the correlation matrix
$\widehat{C}_t = \widehat{D}_t^{-1/2}\, \widehat{\Sigma}_t\, \widehat{D}_t^{-1/2}$,
where $\widehat{D}_t$ holds the diagonal of $\widehat{\Sigma}_t$
\citep{riskmetrics1996}. The clustering map takes the correlation matrix into a
partition of the assets,
\begin{equation}
Q_t \;=\; f\!\left(\widehat{C}_t;\, P\right),
\label{eq:clustermap}
\end{equation}
where $P$ collects the clustering parameters. In our implementation $f$ applies
Ward linkage \citep{ward1963hierarchical} to the distance matrix
$1-\widehat{C}_t$ and cuts the dendrogram at a fixed fraction of its maximal
merge height. The theory below holds $f$ measurable and treats the Ward
specification as one instance.

The clustering is re-estimated at scheduled dates spaced $k$ observation
periods apart. We study the increments
$\Delta\widehat{\Sigma}_t = \widehat{\Sigma}_t - \widehat{\Sigma}_{t-k}$ and the
induced partition changes $\Delta Q_t$, measured by two statistics. The
\emph{churn} is the fraction of assets whose cluster assignment changes between
consecutive scheduled dates, matched through cluster identities, and annualised
to changes per asset-year. The \emph{variation of information} is the metric
$\VI(Q, Q') = H(Q \mid Q') + H(Q' \mid Q)$ on partitions \citep{meila2007comparing}.
Churn is the practice-side statistic and $\VI$ is the theory-side metric. Both
appear in the empirical section so that theory and exhibits speak one language.
The volatility dynamics of returns carry a parameter set $\Theta$, for which we
use GARCH-type specifications \citep{bollerslev1986generalized}.

% ----------------------------------------------------------------------------
\section{The induced partition process and its stationary law}
% ----------------------------------------------------------------------------

%% MECHANISM (candidate for owner rewrite): the point of this section in one
%% sentence: the partition inherits a stationary law from the estimator for
%% free, so the scientific question is not whether the cluster distribution
%% exists but where it concentrates.
The recursion in equation~\eqref{eq:ewma} is a contraction, so the estimated
covariance forgets its initial condition geometrically. Existence of a
stationary cluster distribution is then an inheritance property, not a deep
fact. We record it as a lemma.

\begin{lemma}[Stationarity of the partition process]
\label{lem:stationarity}
Let $\{r_t\}$ be strictly stationary and ergodic with
$\E \log^{+} \| r_1 r_1^{\top} \| < \infty$. Then the recursion
\eqref{eq:ewma} admits the almost-surely convergent stationary solution
\begin{equation}
\widehat{\Sigma}_t \;=\; (1-\lambda) \sum_{j=0}^{\infty} \lambda^{j}\,
r_{t-j} r_{t-j}^{\top},
\label{eq:stationarysol}
\end{equation}
and $\{\widehat{\Sigma}_t\}$ is strictly stationary and ergodic. For any
measurable clustering map $f$, the partition process
$Q_t = f(\widehat{C}_t; P)$ is strictly stationary and ergodic on the finite
lattice of partitions of $d$ assets, and its stationary law is the pushforward
of the stationary law of $\widehat{C}_t$ under $f$.
\end{lemma}

\begin{proof}
The series \eqref{eq:stationarysol} converges almost surely under the
logarithmic moment condition, because $\lambda^j$ decays geometrically while
$\log^{+}\|r_{t-j} r_{t-j}^{\top}\|$ grows sub-linearly along stationary
sequences. The recursion is an iterated random function with a constant
contraction factor $\lambda < 1$, so the backward iteration argument of
\citet{diaconis1999iterated} gives uniqueness of the stationary law and
ergodicity. Measurable functions of stationary ergodic processes are stationary
and ergodic, which covers $Q_t$ and any partition statistic.
\end{proof}

\begin{remark}[Integrated GARCH]
\label{rem:igarch}
The EWMA filter is the integrated GARCH(1,1) filter \citep{riskmetrics1996}.
Under an integrated GARCH data-generating process, returns remain strictly
stationary while their unconditional variance is infinite
\citep{nelson1990stationarity}. Lemma~\ref{lem:stationarity} still applies,
because it requires only the logarithmic moment. Moment-based functionals of
the partition process are the objects that need care in this regime, and the
variance-based results below assume finite fourth moments.
\end{remark}

\begin{remark}[Where the stationary law concentrates]
\label{rem:concentration}
The content of the stationary law is its concentration, not its existence.
Under a block-structured population correlation matrix, exact recovery of the
block partition from $n$ observations requires the within-versus-between
separation to exceed a threshold of order $\sqrt{\log d / n}$
\citep{bunea2020model}. The EWMA estimator carries an effective sample size of
order $N$ (Proposition~\ref{prop:increments}), which does not grow with $d$,
and the random-matrix regime of exponentially weighted estimators keeps the
ratio $d/N$ finite \citep{pafka2004exponential}. The stationary law of $Q_t$
therefore concentrates on the population partition when separation clears the
noise floor and stays diffuse below it. We use this dichotomy as a working
frame and defer its full characterisation, including the removal of the market
mode before clustering \citep{laloux1999noise, lopezdeprado2020machine}, to a
companion study.
\end{remark}

% ----------------------------------------------------------------------------
\section{The law of covariance and correlation increments}
% ----------------------------------------------------------------------------

%% MECHANISM (candidate for owner rewrite): between two monthly estimation
%% dates the EWMA replaces only a few percent of its weight with new data, so
%% the correlation innovations are small and asymptotically Gaussian, and the
%% relevant scale is the estimator's own sampling noise.
The one-step increment of the estimator is exact and already informative. From
equation~\eqref{eq:ewma},
\begin{equation}
\widehat{\Sigma}_t - \widehat{\Sigma}_{t-1}
\;=\; (1-\lambda)\left( r_t r_t^{\top} - \widehat{\Sigma}_{t-1} \right),
\label{eq:increment}
\end{equation}
so the estimator mean-reverts toward the conditional second moment at rate
$1-\lambda$, with an innovation driven by the outer product $r_t r_t^{\top}$.
The next proposition collects the fluctuation results that calibrate all later
formulas.

\begin{proposition}[Increment law and noise floor]
\label{prop:increments}
Let returns be i.i.d.\ with finite fourth moments, let
$V = \operatorname{Var}\!\left(\vech(r_1 r_1^{\top})\right)$, and let
$\lambda = 1 - 2/(N+1)$. Then the following hold.
\begin{enumerate}
\item[(i)] The stationary process $\vech(\widehat{\Sigma}_t)$ has
autocovariance function $\gamma(h) = \lambda^{|h|} \gamma(0)$ with
\begin{equation}
\gamma(0) \;=\; \frac{1-\lambda}{1+\lambda}\, V \;=\; \frac{V}{N},
\label{eq:statvar}
\end{equation}
so the EWMA carries the effective sample size $N$ exactly.
\item[(ii)] The increment over $k$ periods satisfies
\begin{equation}
\operatorname{Var}\!\left( \vech(\widehat{\Sigma}_t)
- \vech(\widehat{\Sigma}_{t-k}) \right)
\;=\; 2\left(1 - \lambda^{k}\right) \gamma(0).
\label{eq:kstep}
\end{equation}
\item[(iii)] For elliptically distributed returns with kurtosis parameter
$\kappa$, the delta method applied to the correlation coefficient gives the
stationary standard error
\begin{equation}
\operatorname{se}\!\left(\widehat{\rho}_{ij}\right)
\;\approx\; \sqrt{1+\kappa}\; \frac{1-\rho_{ij}^2}{\sqrt{N}},
\label{eq:noisefloor}
\end{equation}
with $\kappa = 0$ in the Gaussian case
\citep{neudecker1990asymptotic}.
\item[(iv)] For a GARCH(1,1) data-generating process with Gaussian innovations
and parameters $\Theta = (\alpha, \beta)$ satisfying
$(\alpha+\beta)^2 + 2\alpha^2 < 1$, the unconditional excess-kurtosis parameter
is
\begin{equation}
\kappa_{\Theta} \;=\; \frac{2\alpha^2}{1 - (\alpha+\beta)^2 - 2\alpha^2},
\label{eq:garchkappa}
\end{equation}
so the volatility dynamics enter the noise floor \eqref{eq:noisefloor} through
one closed-form multiplier \citep{bollerslev1986generalized}.
\end{enumerate}
\end{proposition}

\begin{proof}
Part (i): $\vech(\widehat{\Sigma}_t)$ is a linear filter of the i.i.d.\
sequence $\vech(r_t r_t^{\top})$ with coefficients $(1-\lambda)\lambda^{j}$, so
$\gamma(0) = (1-\lambda)^2 V / (1-\lambda^2) = V(1-\lambda)/(1+\lambda)$, and
substituting $\lambda = 1 - 2/(N+1)$ gives $V/N$. The autocovariance follows
from the same filter representation. Part (ii) is the standard identity for
stationary AR(1)-type autocovariances. Part (iii) applies the delta method for
correlations to the fourth-moment structure of elliptical laws, for which the
asymptotic variance of $\widehat{\rho}$ is $(1+\kappa)(1-\rho^2)^2$ per
effective observation. Part (iv) rearranges the unconditional kurtosis of the
Gaussian GARCH(1,1), $K = 3\,(1-(\alpha+\beta)^2)/(1-(\alpha+\beta)^2-2\alpha^2)$,
through $K = 3(1+\kappa_{\Theta})$.
\end{proof}

The distance entries $d_{ij} = 1 - \widehat{\rho}_{ij}$ inherit the standard
error \eqref{eq:noisefloor}. Combining parts (ii) and (iii) gives the
per-transition distance noise at estimation step $k$,
\begin{equation}
\sigma_d(k) \;=\; \sqrt{2\left(1-\lambda^{k}\right)}\;
\sqrt{1+\kappa}\; \frac{1-\bar{\rho}^2}{\sqrt{N}},
\label{eq:transitionnoise}
\end{equation}
evaluated at a representative within-cluster correlation $\bar{\rho}$. Under
misspecification, the same objects follow from the asymptotic filtering theory
of EWMA-type variance estimators \citep{nelson1992filtering,
nelson1994asymptotic}.
% [TODO: check the multivariate filtering reference (J. Econometrics 71, 1996)
% before citing it here.]

Equation \eqref{eq:transitionnoise} quantifies the paper's motivating
arithmetic. At weekly frequency with span $N=156$ and monthly estimation
($k \approx 4.3$), the factor $1-\lambda^{k}$ equals $0.054$: the estimator
replaces about five percent of its weight between estimation dates. The
baseline empirical churn nevertheless reassigns roughly one quarter of assets
per month. A five percent information innovation that produces a twenty-five
percent membership reassignment identifies the dendrogram cut as a noise
amplifier, and the next section locates the amplification at the cluster
boundaries.

% ----------------------------------------------------------------------------
\section{Partition changes: a boundary-flip approximation}
% ----------------------------------------------------------------------------

%% MECHANISM (candidate for owner rewrite): the clustering map is flat almost
%% everywhere and jumps at assignment boundaries, so all churn comes from
%% assets whose margin is small relative to the estimator noise; smoothing is
%% hysteresis that widens the margin.
The clustering map is piecewise constant in the correlation matrix, so no delta
method applies to $\Delta Q_t$: the map is non-differentiable exactly at the
configurations that generate churn. The tractable route fixes a separated
population structure and studies which assets sit close to an assignment
boundary.

\begin{assumption}[Separated block structure]
\label{ass:block}
The population correlation matrix $C^{*}$ has a block structure with $G$
groups, in the sense of the $G$-block covariance models of
\citet{bunea2020model}, and the population dendrogram cut recovers the blocks.
For asset $i$, the population margin
\begin{equation}
m_i \;=\; \bar{d}\left(i, \mathcal{C}_{\mathrm{alt}(i)}\right)
- \bar{d}\left(i, \mathcal{C}(i)\right)
\;>\; 0
\label{eq:margin}
\end{equation}
is the gap between the mean distance from $i$ to its nearest alternative block
and the mean distance to its own block.
\end{assumption}

The margin in \eqref{eq:margin} uses mean distances, which is the
average-linkage functional. Ward linkage does not admit a local margin of this
form, because its merge criterion is a global functional of the tree. We
therefore prove the flip approximation for the flat cut under
Assumption~\ref{ass:block} and verify it for Ward by simulation. This division
is honest rather than convenient: among linkage methods only single linkage is
stable in the Gromov--Hausdorff sense \citep{carlsson2010characterization}, so
a distributional theory for Ward partitions must run through a generative
model, not through generic metric stability.

\begin{proposition}[Boundary-flip approximation]
\label{prop:flip}
Under Assumption~\ref{ass:block}, let the estimated margin statistic
$\widehat{m}_{i,t}$ be asymptotically Gaussian around $m_i$ with standard
deviation $\sigma_i$ proportional to the per-transition distance noise
$\sigma_d(k)$ of equation \eqref{eq:transitionnoise}, with a proportionality
constant absorbing cluster sizes and the correlation of estimation errors.
Under a hysteresis discount $\delta \geq 0$ that retains asset $i$ in its
incumbent cluster unless the estimated margin moves against it by more than
$\delta$, the per-transition flip probability satisfies
\begin{equation}
\Prob\!\left(\text{flip}_i\right)
\;\approx\; \Phi\!\left( - \frac{m_i + \delta}{\sqrt{2}\,\sigma_i} \right),
\label{eq:flipprob}
\end{equation}
and the expected churn per asset-year is
\begin{equation}
\E\!\left[\text{churn}\right]
\;\approx\; \frac{A}{d} \sum_{i=1}^{d}
\Phi\!\left( - \frac{m_i + \delta}{\sqrt{2}\,\sigma_i} \right),
\label{eq:expectedchurn}
\end{equation}
where $A$ is the number of estimation dates per year.
\end{proposition}

\begin{proof}[Proof sketch]
A flip requires the noisy margin to cross zero against the hysteresis band,
which is a Gaussian tail event for the difference of two estimated mean
distances, whence the factor $\sqrt{2}$. Summing marginal flip probabilities
gives the expected count without any independence requirement, because
expectation is linear. The Gaussian approximation for $\widehat{m}_{i,t}$
follows from Proposition~\ref{prop:increments} and the delta method under
Assumption~\ref{ass:block}, where merge heights concentrate and the population
tree is deterministic.
\end{proof}

%% MECHANISM (candidate for owner rewrite): the design principle in one line:
%% a membership change should require statistical evidence that exceeds the
%% estimator's own sampling noise, which is the same gate logic as the
%% sign-constraint gate of our companion estimator.
Equation \eqref{eq:expectedchurn} turns the smoothing parameter into a designed
quantity. Two calibrations follow, and they differ in an observable way.

\begin{corollary}[Noise-floor calibration of the partition bonus]
\label{cor:delta}
Fix a false-flip budget $\alpha$ for zero-margin assets, with Gaussian quantile
$z_{1-\alpha}$. The level calibration sets the bonus against the stationary
noise of the distance level,
\begin{equation}
\delta^{*}_{\mathrm{lvl}}
\;=\; z_{1-\alpha}\, \sqrt{1+\kappa}\; \frac{1-\bar{\rho}^{\,2}}{\sqrt{N}},
\label{eq:deltalevel}
\end{equation}
and the innovation calibration sets it against the per-transition noise,
\begin{equation}
\delta^{*}_{\mathrm{inn}}
\;=\; \sqrt{2\left(1-\lambda^{k}\right)}\;\, \delta^{*}_{\mathrm{lvl}}.
\label{eq:deltainnovation}
\end{equation}
The level form does not depend on the estimation step $k$, while the
innovation form shrinks as estimation dates move closer together. The
dependence on $k$ therefore discriminates between the two calibrations
empirically.
\end{corollary}

At the parameters of our weekly implementation ($N=156$, $k \approx 4.3$,
$\bar{\rho} \approx 0.5$, Gaussian $\kappa = 0$, $z_{1-\alpha} = 1$), the level
form gives $\delta^{*}_{\mathrm{lvl}} \approx 0.060$ and the innovation form
gives $\delta^{*}_{\mathrm{inn}} \approx 0.020$. The value selected by our
parameter sweeps, $\delta = 0.05$, lies between the two. We regard this
bracketing as encouraging rather than conclusive, and
Prediction~\ref{pred:frequency} below states the test that separates the
calibrations.

% ----------------------------------------------------------------------------
\section{The smoothing design problem}
% ----------------------------------------------------------------------------

Minimising partition variability alone is degenerate: any constant partition
achieves zero churn, so an unconstrained stability objective selects a frozen
or trivial clustering. This degeneracy is the central caveat of the clustering
stability literature, where stability measures boundary mass rather than
correctness \citep{bendavid2006sober, vonluxburg2010clustering}. We therefore
state the design problem as constrained:
\begin{equation}
\min_{P} \;\; \E\!\left[\text{churn}(P)\right]
\qquad \text{subject to} \qquad
\mathcal{F}(P) \;\geq\; \mathcal{F}_0,
\label{eq:designproblem}
\end{equation}
where $\mathcal{F}$ is a fidelity functional. In the theory, $\mathcal{F}$ is
the probability of recovering the population partition under
Assumption~\ref{ass:block}. In the implementation, $\mathcal{F}$ is the
adjusted Rand index band against the unsmoothed baseline together with a
resolution guard on the cluster-count distribution.

\begin{remark}[Relation to evolutionary clustering]
\label{rem:evolutionary}
The Lagrangian form of \eqref{eq:designproblem} is the snapshot-cost plus
temporal-cost objective of evolutionary clustering
\citep{chakrabarti2006evolutionary}, and similarity smoothing corresponds to
the adaptive forgetting estimators of that line \citep{xu2014adaptive}. Our
contribution at this point is not the objective but its calibration: the
temporal weight is derived from the estimator's noise floor through
Corollary~\ref{cor:delta} rather than tuned. Sample-length requirements for
correlation-based clustering in finance are studied by
\citet{marti2016clustering}, whose rates concern the time dimension of the
same trade-off.
\end{remark}

% ----------------------------------------------------------------------------
\section{Empirical implications}
% ----------------------------------------------------------------------------

The theory is falsifiable on the three data sets of the empirical section
(United States equities with a six-factor model at weekly frequency, global
futures and a multi-asset fund universe with a custom factor model at weekly,
monthly and quarterly frequencies). We state the implications as numbered
predictions. Each names its test statistic, and the empirical section reports
each with a bootstrap confidence interval or a permutation null.
% Internal mapping (not for the manuscript): P1-P4, P6 -> stage E3;
% P5 -> stage E3 (risk-model invariance); P7 -> stage E5; inference -> E6.

\begin{prediction}[Churn concentrates at small margins]
\label{pred:margins}
Realised membership flips concentrate on assets whose estimated margins are
small relative to $\sigma_d(k)$, and the predicted churn
\eqref{eq:expectedchurn} matches realised churn across smoothing
configurations. Test: per-universe margin histograms, flip rates by margin
decile, and the cross-configuration correlation between predicted and realised
churn.
\end{prediction}

\begin{prediction}[The churn-fidelity frontier turns at the noise floor]
\label{pred:knee}
On the frontier traced by the bonus parameter $\delta$, churn falls steeply up
to the noise-floor calibration of Corollary~\ref{cor:delta} and fidelity
deteriorates beyond it. Test: sweep $\delta$ on a grid, overlay
$\delta^{*}_{\mathrm{lvl}}$ and $\delta^{*}_{\mathrm{inn}}$ computed from each
universe's span, frequency and kurtosis, and locate the frontier knee.
\end{prediction}

\begin{prediction}[Kurtosis raises the noise floor]
\label{pred:kurtosis}
Universes with fatter tails carry proportionally higher churn at a fixed
$\delta$, with the ordering given by $\sqrt{1+\kappa}$ estimated from the
data. Test: cross-universe comparison of realised churn against the
kurtosis-adjusted prediction, with $\kappa$ estimated from pooled returns and,
under the GARCH specification, from equation \eqref{eq:garchkappa}.
\end{prediction}

\begin{prediction}[Frequency scaling discriminates the calibrations]
\label{pred:frequency}
Per-transition churn scales as $\sqrt{1-\lambda^{k}}$ when estimation dates
move from monthly to quarterly on the same universe. The optimal bonus
$\delta^{*}$ moves with $k$ under the innovation calibration
\eqref{eq:deltainnovation} and stays fixed under the level calibration
\eqref{eq:deltalevel}. Test: re-score cached estimations at the coarser
schedule and compare both scalings against their predictions.
\end{prediction}

\begin{prediction}[Stabilisation is free at the risk-model level]
\label{pred:invariance}
At production regularisation, the fitted covariance is insensitive to the
partition, so smoothing changes the risk model by an order of magnitude less
than it changes the partition. Test: relative Frobenius distance and maximal
entry change of the fitted covariance, and the residual-diagonality
diagnostics, between smoothed and baseline runs on all three universes.
\end{prediction}

\begin{prediction}[Ergodic partition statistics]
\label{pred:ergodicity}
Partition summary statistics (cluster count, size entropy, consecutive-date
$\ARI$ and $\VI$) have stable subsample averages over the sample, consistent
with the stationarity and ergodicity of Lemma~\ref{lem:stationarity}. Test:
subsample means across sample halves and thirds against bootstrap bands, with
crisis windows reported separately.
\end{prediction}

\begin{prediction}[Turnover attribution in the momentum overlay]
\label{pred:turnover}
In a cluster-relative momentum portfolio, reassignment-driven turnover falls
monotonically in the smoothing strength while signal-driven turnover is
invariant, and net-of-cost performance is non-decreasing within the fidelity
band. Test: turnover decomposition against the prior-date partition
counterfactual, at transaction costs of 10, 20 and 50 basis points on the
three universes respectively, with block-bootstrap confidence intervals on the
turnover and net performance deltas.
\end{prediction}

%% MECHANISM (candidate for owner rewrite): the implications paragraph after
%% the main result, per the revision criteria.
Two consequences follow if the predictions hold. For practice, the smoothing
parameter stops being a tuning choice: a risk platform can set $\delta$ from
the span, the estimation frequency and a tail estimate, and defend the setting
as a statistical-evidence threshold rather than a preference. For the
literature, the boundary-flip frame shifts the object of study from the
clustering algorithm to the interaction between estimator noise and partition
geometry, which is where the empirically observed churn actually originates.
If Prediction~\ref{pred:frequency} rejects both calibrations, the hysteresis
model in Proposition~\ref{prop:flip} is the component to revisit first,
because the remaining results do not depend on it.

\bibliographystyle{plainnat}
\bibliography{cluster_lineage_refs}

\end{document}
